% Chapter 2 exercises. The answers are in answers/barcs_answers_ch2.pdf.
\section{Exercises for Chapter 2}

These exercises continue the screen of the Chapter~1 exercises (guide g7).
Work on paper and show each calculation. The answers are in a separate
answer sheet.

The tags tell you what to do: \tagcalc{} work out a number,
\tagexplain{} say why in your own words, \tagdecide{} choose and give a
reason.

\begin{center}
\fbox{\begin{minipage}{0.92\linewidth}
\textbf{Values to use.} Expected proportions: $\mu = 0.00015$ (control) and
$\mu = 0.000075$ (drug). Biological noise: $\rho = 0.000001$. Effect:
$\beta_1 = -0.69$. Use $\ln$ for the natural logarithm.
\end{minipage}}
\end{center}

\subsection*{Part A: Fitting}
\begin{exlist}[series=ex2]
\item \tagcalc{} \tagexplain{} A control library has $K/N = 0.00018$. The
current prediction is $\mu = 0.00015$, so $\eta = \ln(0.00015/0.99985) =
-8.8047$.
\begin{parts}
\item Calculate the slope $1 \div \bigl(\mu \times (1-\mu)\bigr)$.
\item Calculate the step $(K/N - \mu) \times \text{slope}$ and the working
  value $v = \eta + \text{step}$.
\item Calculate the real log-odds $\ln(0.00018/0.99982)$ and compare.
\item Why does BARCS use the straight-line working value and not the real
  log-odds? Think of a guide with $K = 0$.
\end{parts}

\item \tagcalc{} \tagexplain{} Two control libraries have $N = 1{,}000{,}000$
and $N = 2{,}000{,}000$. Use $\mu = 0.00015$.
\begin{parts}
\item Calculate the weight
  $w = N \times \mu \times (1-\mu) \div \bigl(1 + (N-1) \times \rho\bigr)$
  of each library.
\item The second library has twice the reads. Does it have twice the weight?
  Explain why not.
\item What would the two weights be if $\rho = 0$?
\end{parts}

\item \tagcalc{} \tagexplain{} Two libraries have the working values
$-8.60$ (weight 75) and $-8.90$ (weight 100).
\begin{parts}
\item Calculate the weighted average.
\item Calculate its standard error, $\sqrt{1 \div (75 + 100)}$.
\item The weighted average is closer to $-8.90$ than to $-8.60$. Why is
  that the right thing to do?
\end{parts}
\end{exlist}

\subsection*{Part B: Is the effect real?}
\begin{exlist}[resume*=ex2]
\item \tagcalc{} \tagexplain{} The control weights add up to
$W_{\text{control}} = 175$. The drug weights add up to $W_{\text{drug}} = 70$.
\begin{parts}
\item Calculate $\text{SE}(\beta_1) = \sqrt{1/W_{\text{control}} +
  1/W_{\text{drug}}}$.
\item Calculate $t = \beta_1 \div \text{SE}$ and the degrees of freedom
  (4 libraries, 2 $\beta$ values).
\item With 2 degrees of freedom, $|t|$ must be larger than 4.30 for
  $p < 0.05$. Is the effect significant?
\item Which group adds more uncertainty? What would you change in the next
  experiment to make the SE smaller?
\end{parts}

\item \tagcalc{} \tagexplain{} Two guides, P and Q, each have 4 libraries and
2 $\beta$ values. At $\rho = 0$, the squared Pearson residuals add up to 5.1
for guide P and to 1.4 for guide Q.
\begin{parts}
\item Calculate the variance inflation (sum $\div$ df) of each guide.
\item For which guide does BARCS use $\rho > 0$? For which does it use
  $\rho = 0$?
\item Why do we compare the sum with 2 and not with 4?
\end{parts}

\item \tagcalc{} \tagexplain{} Guide P has its own inflation 2.55 with
2 degrees of freedom. In the screen, the typical inflation is 1.5 with
$d_0 = 10$.
\begin{parts}
\item Calculate the moderated inflation.
\item Calculate the factor for the SE, $\sqrt{\text{moderated} \div
  \text{own}}$, and the new degrees of freedom.
\item Why is it reasonable to borrow information from other guides? When
  could it be a bad idea?
\end{parts}
\end{exlist}

\subsection*{Part C: From guides to a list of genes}
\begin{exlist}[resume*=ex2]
\item \tagcalc{} \tagdecide{} The screen has 300 non-targeting control
guides. 45 of them get $p < 0.05$.
\begin{parts}
\item How many would you expect if the noise model were correct? What does
  the real number tell you?
\item 95\% of the control guides have $|t|$ below 2.6. The $t$ table says
  1.96. Calculate the scale factor.
\item A gene has $t = 3.0$. Calculate its new $t$. Is it still beyond 1.96?
\end{parts}

\item \tagcalc{} \tagexplain{} A gene has four guides with the z scores 2.2,
1.8, $-0.3$ and 2.5.
\begin{parts}
\item Calculate $Z_{\text{gene}}$ with the Stouffer method (no shared
  noise).
\item Calculate $Z_{\text{gene}}$ again with the guide correlation
  $r = 0.2$.
\item Why does one guide with a negative $z$ lower the result? Why does
  shared noise lower the result?
\end{parts}

\item \tagcalc{} \tagdecide{} A screen tests 10{,}000 genes. 40 genes have
$p \le 0.0005$.
\begin{parts}
\item If no gene had an effect, about how many genes would have
  $p \le 0.0005$ by chance?
\item Estimate the false discovery rate of this list of 40 genes.
\item You want a list with an FDR of about 0.05. Should you use a smaller or
  a larger p-value limit? Explain.
\end{parts}

\item \tagdecide{} \tagexplain{} A collaborator says: ``We will read 20
million reads per library. With so many reads, 2 replicates per group are
enough.'' Use $\rho = 0.000001$ and the degrees of freedom to answer.
(Hint: calculate the extra factor $1 + (N-1) \times \rho$ for
$N = 20{,}000{,}000$.)

\item \tagcalc{} \tagexplain{} \tagdecide{} (Step~13) In the control check of
Exercise~7, the control guides still have too many extreme values after
scaling. You try the empirical null.
\begin{parts}
\item A gene has 3 guides with the z scores 2.0, 2.6 and 2.2. Calculate the
  average z score.
\item BARCS draws 1{,}999 pseudo-genes of 3 control guides. 9 of them have
  an average of your answer to (a) or more. Calculate the p-value.
\item The guide correlation is $r = 0.1$. By what factor does BARCS widen
  each pseudo-gene? Why is this needed?
\item The screen has only 150 control guides and many genes with one guide.
  What is the smallest p-value a one-guide gene can get? What would you
  tell a student who sees ``no hits'' from the empirical null?
\end{parts}
\end{exlist}
